\begin{definition}[Literal directed transports after a double exchange]%
    \label{def:peps_torus_three_plaquette_double_exchange_transports}
    \lean{TNLean.PEPS.torusThreePlaquetteDoubleExchangeTransports}
    \leanok
    \uses{def:peps_torus_three_plaquette_flux_assignment,
        def:peps_regular_three_plaquette_double_exchange}
    For directed transports $t=(t_0,t_1,t_2)$, define
    \begin{align}
        t'=(t_0,\ t_2t_0^{-1}t_1t_2^{-1}t_0,\ t_2).
    \end{align}
    If $C(t)=(a,b,c)$ denotes the three plaquette holonomies, then
    $t'=C^{-1}H^2C(t)$, with $H(a,b)=(aba^{-1},a)$ acting on the first
    two plaquettes. The same ordered-bond orientation correction applies
    to both literal assignments.
\end{definition}

\begin{theorem}[Actual holonomies of the literal double-exchange output]%
    \label{thm:peps_torus_three_plaquette_double_exchange_holonomy}
    \lean{TNLean.PEPS.regularWalkHolonomy_torusThreePlaquetteDoubleExchange}
    \leanok
    \uses{def:peps_torus_three_plaquette_double_exchange_transports,
        thm:peps_torus_three_plaquette_flux_holonomy}
    At every translated position, the literal output has plaquette
    holonomies
    \begin{align}
        (a',b',c')=((ab)a(ab)^{-1},\ aba^{-1},\ c).
    \end{align}
    The outer middle-right walk has holonomy $b'c'=aba^{-1}c$.
    In particular, at $(a,b,c)=(g,h,h^{-1})$ this surrounding holonomy
    is $ghg^{-1}h^{-1}$.
    The original period bounds, at least five horizontally and three
    vertically, suffice also at both seams.
\end{theorem}

\begin{proof}\leanok
    Substitute $t'$ into the actual directed-walk formulas.
    Group cancellation gives the displayed plaquette and outer holonomies.
\end{proof}

\begin{theorem}[One original-spin double exchange on the translated strip]%
    \label{thm:peps_torus_three_plaquette_double_exchange_physical}
    \lean{TNLean.PEPS.IsGIsometric.exists_unitary_torusThreePlaquetteDoubleExchange}
    \leanok
    \uses{thm:peps_three_plaquette_geometry,
        thm:peps_torus_three_plaquette_double_exchange_holonomy,
        thm:peps_regular_three_plaquette_double_exchange_physical}
    Let $G$ be finite, and suppose the native four-leg tensor is
    $G$-isometric for the regular torus-leg representation.
    For every position $v$, there is one unitary $W$ on the eight original
    spins of $R_v$, chosen before every directed transport triple $t$,
    boundary configuration $\theta$, and common exterior and crossing
    assignment $u$, such that
    \begin{align}
        W\Psi_{R_v}(t,\theta)         & =\Psi_{R_v}(t',\theta), \\
        (W\otimes I_{R_v^c})\Psi(t,u) & =\Psi(t',u).
    \end{align}
    Both are actual original-spin contractions with the literal assignments
    just defined; the identity extension is unitary.
    The region, spanning tree, root, three chords and their orientation are
    derived from the translated strip geometry.
    This is an auxiliary finite realization of the double-exchange
    algebra in \cite[braiding passage]{Schuch2010PEPS}; it does not identify
    the operation with the prescribed physical string crossing or assert
    parent-Hamiltonian membership.
\end{theorem}

\begin{proof}\leanok
    The native four-leg hypothesis implies incident-edge $G$-isometry.
    Apply the common physical double-exchange theorem to the derived tree,
    root and three chords, using the common seam-orientation flag.
    The actual literal-to-tree-cycle identity identifies the input columns.
    The sparse residual calculation identifies the output with the literal
    transport triple $t'$. Both the local and global identities therefore
    use the same supplied unitary.
\end{proof}
